% 第 8 章 应用前景与预期效果 (目标 1.5 页)
% 对应官方要求: 预期效果
\section{应用前景与预期效果}

\subsection{落地形态}

本系统代码随报告提交,核心模块含结构化规则引擎、结构化评审层、真实环境接入代理、攻击与白样本评测集,均提供运行说明。第三方团队可在本地智能体上重现全部实验。部署支持单机代理与内网网关两种形态,适配不同规模的政企单位。

\subsection{推广路径}

规则层与评审层均与具体业务技能解耦。规则通过新增编号与谓词即可覆盖新的高风险工具或参数模式,评审通过修改指令即可调整评估维度。两层均不依赖特定模型,可在多模型环境部署。政务办公之外,金融、医疗、能源等数据敏感行业的智能体防护可直接复用本架构,按行业补充规则库即可。

\subsection{对治理水平的价值}

本系统给出从关键词到结构化规则到结构化评审的三阶段演进路径,可作为政企单位建设智能体安全防护体系的参考模板。审计溯源设计满足合规留痕要求,评测数据集与评测方法可复用为行业评测基准的种子。

\subsection{预期效果}

按当前实验结果外推,系统部署后可将去标记化业务话术诱导类攻击的绕过率从 86.4\% 压到 0\%,正常业务操作零误伤,单次调用增加延迟 1.8 秒以内。随着规则库与评审指令在行业部署中持续迭代,覆盖率与误报控制仍有进一步提升空间。
